%-------------------------------------------
% macros.tex
% The resume/CV switch and the commands every
% section uses to lay out entries.
%
% Based on Jake's Resume by Jake Gutierrez
% (github.com/jakegut/resume, MIT), itself
% based on github.com/sb2nov/resume (MIT).
%-------------------------------------------

%----------RESUME OR CV----------
% cv.tex defines \buildcv before loading this file; any other root
% (resume.tex) builds the one-page resume.
\newif\ifcv
\ifdefined\buildcv \cvtrue \fi

% Write each entry once and mark only what differs between the versions:
%   \cvonly{...}             CV only
%   \resumeonly{...}         resume only
%   \cvresume{CV}{resume}    different text in each
% Arguments can't hold verbatim text; escape any % or # in a URL.
\newcommand{\cvonly}[1]{\ifcv\ignorespaces#1\fi\ignorespaces}
\newcommand{\resumeonly}[1]{\ifcv\else\ignorespaces#1\fi\ignorespaces}
\newcommand{\cvresume}[2]{%
  \ifcv\ignorespaces#1\else\ignorespaces#2\fi\ignorespaces}

%----------ENTRY COMMANDS----------
% \unskip drops the trailing space of #1; without it, a bullet that exactly
% fills its last line breaks at that space and adds an empty line.
\newcommand{\resumeItem}[1]{
  \item\small{
    {#1\unskip\vspace{-2pt}}
  }
}

% One full-width row: #1 flush left, #2 flush right. Every heading below
% is built from it. \nopagebreak keeps each row on the same page as what
% follows it: the next row of the heading, or its first bullet.
\newcommand{\resumeRow}[2]{%
  \begin{tabular*}{0.97\textwidth}[t]{l@{\extracolsep{\fill}}r}
    #1 & #2 \\
  \end{tabular*}\nopagebreak\vspace{-7pt}%
}

% Bold organization (employer or university) and its location.
\newcommand{\resumeOrg}[2]{%
  \vspace{-2pt}\item
  \resumeRow{\textbf{#1}}{#2}%
}

% Italic title (job title or degree) and its dates. Repeat it for
% promotions; bullets may follow any title.
\newcommand{\resumeTitle}[2]{%
  \vspace{-2pt}\item
  \resumeRow{\textit{\small#1}}{\textit{\small #2}}%
}

\newcommand{\resumeProjectHeading}[2]{%
  \item
  \resumeRow{\small#1}{#2}%
}

\newcommand{\resumeSubItem}[1]{
  \resumeItem{#1}\vspace{-4pt}
}

\renewcommand\labelitemii{
  $\vcenter{\hbox{\tiny$\bullet$}}$
}

\newcommand{\resumeSubHeadingListStart}{
  \begin{itemize}[leftmargin=0.15in, label={}]
}
\newcommand{\resumeSubHeadingListEnd}{\end{itemize}}
\newcommand{\resumeItemListStart}{\begin{itemize}}
\newcommand{\resumeItemListEnd}{
  \end{itemize}\vspace{-5pt}
}

%----------PUBLICATIONS----------
% \pubitem{key}{citation} numbers an entry in the Publications list;
% \pubcite{key} prints its number, e.g. [2], anywhere else.
\newcounter{pubcounter}
\newcommand{\pubitem}[2]{
  \refstepcounter{pubcounter}
  \label{pub:#1}
  \item[{\small[\thepubcounter]}] #2
}
\newcommand{\pubcite}[1]{[\ref{pub:#1}]}
